% EvoClaw: A Self-Evolving Multi-Agent Framework
% with Genome-Driven Runtime Adaptation, Tiered Memory, 
% and Intelligent Model Routing
%
% Extended 50+ Page Version — arXiv cs.AI / cs.MA submission
% Targeting NeurIPS 2024 format
%
% Comprehensive coverage of:
% - Genome-driven evolution
% - Tiered memory architecture
% - Intelligent model routing
% - Model health monitoring
% - Agentic tool loops
% - Self-governance protocols
% - Blockchain integration
% - Multi-tier deployment

\documentclass{article}
\usepackage[preprint]{neurips_2024}
\usepackage[utf8]{inputenc}
\usepackage[T1]{fontenc}
\usepackage{hyperref}
\usepackage{url}
\usepackage{booktabs}
\usepackage{amsfonts}
\usepackage{amsmath}
\usepackage{amssymb}
\usepackage{amsthm}
\usepackage{algorithm}
\usepackage{algpseudocode}
\usepackage{tikz}
\usepackage{pgfplots}
\usepackage{listings}
\usepackage{xcolor}
\usepackage{subcaption}
\usepackage{multirow}
\usepackage{enumitem}
\usepackage{wrapfig}
\usepackage{longtable}
\usepackage{array}

% TikZ libraries
\usetikzlibrary{arrows.meta, positioning, shapes.geometric, shapes.multipart,
  fit, calc, decorations.pathrepeating, backgrounds, shadows.blur,
  patterns, chains, matrix, trees}

% PGFPlots settings
\pgfplotsset{compat=1.18}

% Color palette
\definecolor{coreblue}{HTML}{2563EB}
\definecolor{coreteal}{HTML}{0D9488}
\definecolor{coreorange}{HTML}{EA580C}
\definecolor{corepurple}{HTML}{7C3AED}
\definecolor{coregray}{HTML}{64748B}
\definecolor{coregreen}{HTML}{16A34A}
\definecolor{corered}{HTML}{DC2626}
\definecolor{lightbg}{HTML}{F8FAFC}
\definecolor{edgebg}{HTML}{ECFDF5}
\definecolor{serverbg}{HTML}{EFF6FF}
\definecolor{cloudbg}{HTML}{FEF3C7}

% Listing styles
\lstdefinestyle{rustcode}{
  language=C,
  basicstyle=\ttfamily\footnotesize,
  keywordstyle=\color{coreblue}\bfseries,
  commentstyle=\color{coregray}\itshape,
  stringstyle=\color{coregreen},
  breaklines=true,
  frame=single,
  rulecolor=\color{coregray!30},
  backgroundcolor=\color{lightbg},
  numbers=left,
  numberstyle=\tiny\color{coregray},
  numbersep=5pt,
  tabsize=4,
  showstringspaces=false,
  morekeywords={fn, let, mut, pub, struct, impl, trait, async, await, Result, Self, String, Vec, Box, dyn, use, mod, crate, where, type, enum, match, Ok, Err, Some, None},
}

\lstdefinestyle{tomlcode}{
  basicstyle=\ttfamily\footnotesize,
  keywordstyle=\color{corepurple}\bfseries,
  commentstyle=\color{coregray}\itshape,
  stringstyle=\color{coregreen},
  breaklines=true,
  frame=single,
  rulecolor=\color{coregray!30},
  backgroundcolor=\color{lightbg},
  numbers=none,
  tabsize=4,
  showstringspaces=false,
}

\lstdefinestyle{gocode}{
  language=C,
  basicstyle=\ttfamily\footnotesize,
  keywordstyle=\color{coreteal}\bfseries,
  commentstyle=\color{coregray}\itshape,
  stringstyle=\color{coregreen},
  breaklines=true,
  frame=single,
  rulecolor=\color{coregray!30},
  backgroundcolor=\color{lightbg},
  numbers=left,
  numberstyle=\tiny\color{coregray},
  numbersep=5pt,
  tabsize=4,
  showstringspaces=false,
  morekeywords={func, go, chan, select, defer, range, interface, struct, package, import, type, map, make, return},
}

% Theorem environments
\newtheorem{definition}{Definition}
\newtheorem{theorem}{Theorem}
\newtheorem{lemma}{Lemma}
\newtheorem{property}{Property}
\newtheorem{assumption}{Assumption}

% Custom commands
\newcommand{\evoclaw}{\textsc{EvoClaw}}
\newcommand{\genome}{\texttt{genome.toml}}
\newcommand{\skill}[1]{\textsc{#1}}
\newcommand{\topic}[1]{\texttt{#1}}
\newcommand{\fitnessfn}{\mathcal{F}}
\newcommand{\genomespace}{\mathcal{G}}
\newcommand{\skillset}{\mathcal{S}}
\newcommand{\mutation}{\mathcal{M}}
\newcommand{\R}{\mathbb{R}}

\title{\evoclaw{}: A Self-Evolving Multi-Agent Framework \\with Genome-Driven Runtime Adaptation, \\Tiered Memory, and Intelligent Model Routing}

\author{
  Alex Chen\thanks{Equal contribution. Correspondence: \texttt{bowen31337@outlook.com}} \quad Bowen Li\footnotemark[1] \quad Nicholas Qi\footnotemark[1] \\[0.3em]
  Independent Researchers \\
  \texttt{alex.chen31337@gmail.com} \quad \texttt{bowen31337@outlook.com} \quad \texttt{nicholas68663@gmail.com} \\
  \texttt{github.com/clawinfra/evoclaw}
}

\begin{document}

\maketitle

\begin{abstract}
Current multi-agent systems are fundamentally static: once deployed, an agent's architecture, capabilities, and behavior remain fixed until human engineers manually update them. This limitation prevents agents from adapting to dynamic environments, learning from operational experience, or evolving autonomously over time. We present \evoclaw{}, a self-evolving multi-agent framework where agents modify their own behavior at runtime through a declarative genome specification. \evoclaw{} introduces six core innovations: (1) a \genome{} declarative specification that encodes an agent's capabilities as a composable set of pluggable skills with evolutionary parameters; (2) a \textit{Genome Engine} implemented in Rust that parses, validates, and hot-reloads skill configurations without downtime; (3) a \textit{tiered memory architecture} inspired by human cognition, with hot/warm/cold tiers and tree-structured retrieval achieving O(log n) search; (4) an \textit{intelligent model router} that classifies tasks by complexity and routes to appropriate LLMs, reducing costs by 10--100$\times$; (5) a \textit{model health registry} implementing circuit-breaker patterns for automatic failover; and (6) \textit{self-governance protocols} (WAL, VBR, ADL, VFM) ensuring reliability without human supervision.

We formalize the genome evolution process as a constrained optimization over a skill configuration space and introduce a fitness function that balances task performance, resource utilization, and safety invariants. We demonstrate \evoclaw{}'s practicality by deploying a fully functional agent on a Raspberry Pi~1 (ARMv6, 512MB RAM) running three concurrent skills---system monitoring, GPIO control, and cryptocurrency price tracking---with a binary size of 3.0\,MB and memory footprint of 3.2\,MB RSS. Our Rust implementation achieves 90\%+ test coverage with 383 tests, and our Go orchestrator manages 500+ concurrent agents via MQTT pub/sub with sub-50\,ms message latency. \evoclaw{} opens research directions in self-modifying systems, edge AI, and autonomous agent economies. The complete system is open-source and fully reproducible.\footnote{Code: \url{https://github.com/clawinfra/evoclaw}}
\end{abstract}

% ========================================================================
% TABLE OF CONTENTS (for 50+ page paper)
% ========================================================================

\tableofcontents

%% ========================================================================
%%  SECTION 1: INTRODUCTION
%% ========================================================================

\section{Introduction}
\label{sec:introduction}

The rise of large language models (LLMs) has enabled a new generation of autonomous agents that plan tasks, write code, use tools, and interact with complex environments~\cite{wang2023survey, xi2023rise}. Systems like AutoGPT~\cite{auto_gpt}, BabyAGI~\cite{babyagi}, and Voyager~\cite{wang2023voyager} demonstrate that LLM-powered agents can decompose high-level goals into executable plans and accumulate skills over time. However, these agents suffer from a fundamental architectural limitation: they are \textit{static}. Once deployed, an agent's capabilities---its available tools, behavioral patterns, and response strategies---remain fixed until human engineers manually rewrite and redeploy the codebase~\cite{wooldridge2009introduction}.

This static architecture creates a critical bottleneck in three dimensions:

\textbf{Adaptability.} Agents cannot modify their own skill repertoire in response to environmental changes (e.g., new APIs, shifting user requirements, hardware failures)~\cite{ge2023metagpt}. 

\textbf{Efficiency.} Resource-constrained edge devices cannot afford to run monolithic agent binaries; agents should activate only the skills relevant to their current context~\cite{shi2016edge}.

\textbf{Scalability.} Managing hundreds of agents across heterogeneous deployment tiers requires a unified abstraction for agent identity, capabilities, and evolution---not bespoke configuration for each instance~\cite{guo2024llm_multi_agents}.

We argue that truly autonomous agents should be able to modify themselves~\cite{vonneumann1966self}. Just as biological organisms evolve through natural selection---adapting their phenotype to environmental pressures---software agents should evolve through \textit{directed mutation}: guided by their own operational experience, environmental feedback, or human oversight~\cite{koza1992gp}. The challenge is designing a system that allows safe, verifiable self-modification without sacrificing reliability, security, or performance~\cite{rushby1995critical}.

\subsection{Key Insight: Agents as Genomes}

Our key insight is that an agent's capabilities can be represented as a \textit{genome}---a declarative specification that is separate from the agent's runtime binary~\cite{smith1984reflection}. By externalizing the agent's skill configuration into a human-readable, machine-parseable file (\genome{}), we achieve a clean separation between the \textit{execution engine} (fixed Rust binary) and the \textit{behavioral specification} (mutable TOML configuration)~\cite{meyer1997object}. This separation enables:

\textbf{Hot-reload without recompilation:} Skill configurations can be modified and reloaded at runtime, enabling zero-downtime adaptation~\cite{hewitt1976implementation}.

\textbf{Evolutionary operators:} Standard genetic programming operators (mutation, crossover, selection) can be applied to the genome, enabling automated capability search~\cite{angeline1994genetic}.

\textbf{Version control and rollback:} Genomes are plain-text files that can be tracked in Git, enabling full audit trails and safe rollback on regression~\cite{tan2000git}.

\textbf{Cross-tier portability:} The same genome format works across edge, server, and cloud deployments, with tier-specific resource constraints enforced at load time~\cite{tinyml}.

\subsection{Beyond Static Agents: Six Core Innovations}

While genome-driven evolution is our foundational contribution, we identify and address five additional critical gaps in contemporary agent systems:

\textbf{1. Tiered Memory Architecture.} Current LLM agents suffer from context window bloat, similarity-based retrieval errors (RAG~\cite{lewis2020retrieval}), and linear scaling~\cite{zhong2023token}. We introduce a three-tier memory system (Hot/Warm/Cold) inspired by human cognition~\cite{atkinson1971human} and PageIndex~\cite{vectifyai_pageindex}, achieving O(log n) tree-structured retrieval instead of O(n) vector search. Core memory stays bounded at 5KB regardless of conversation history.

\textbf{2. Intelligent Model Routing.} LLM API costs dominate agent operating expenses~\cite{liu2023llm}. We design a router that classifies tasks by complexity (Simple/Medium/Complex/Reasoning/Critical) using 15 weighted dimensions, routing simple tasks to \$0.20/M token models and critical tasks to \$15/M models. This reduces costs by 10--100$\times$ while maintaining quality via automatic fallback chains~\cite{frugal}.

\textbf{3. Model Health Registry.} Model outages, rate limits, and quota exhaustion cause cascading failures in production systems~\cite{tan2020head}. We implement a circuit-breaker pattern~\cite{ford2000circuit} that tracks per-model success rates, automatically degrading unhealthy models and routing to fallbacks. State persists across restarts for resilience.

\textbf{4. Agentic Tool Loop.} Edge agents have access to 30+ tools (file ops, bash, git, browser) but LLMs cannot execute them without structured schemas~\cite{schick2023toolformer}. We implement multi-turn tool loops where LLMs generate tool calls, edge agents execute them, and results feed back into conversation history, enabling complex multi-tool workflows~\cite{parisi2023gorilla}.

\textbf{5. Self-Governance Protocols.} Autonomous agents need reliability mechanisms without human supervision~\cite{shinn2024reflexion}. We implement four protocols: Write-Ahead Log (WAL) prevents context loss during compaction~\cite{hewitt1976implementation}; Verify Before Reporting (VBR) prevents false completion claims; Anti-Divergence Limit (ADL) prevents persona drift~\cite{park2023generative}; and Value-For-Money (VFM) prevents wasteful spending.

\textbf{6. Blockchain Integration.} Looking ahead, genome-driven agents enable reputation-based coordination via ClawChain~\cite{clawchain}, our agent-native Layer~1 blockchain. Agents register on-chain DIDs~\cite{spencer2015decentralized}, accumulate reputation through verified actions, and participate in governance---creating a trustless agent economy~\cite{buterin2014ethereum}.

\subsection{Contributions}

We present \evoclaw{}, a multi-agent framework where agents evolve through a declarative genome specification with six integrated innovations. Our contributions are:

\begin{enumerate}[leftmargin=*]
    \item \textbf{Genome-Driven Architecture} (\S\ref{sec:architecture}): We introduce \genome{}, a declarative configuration language that specifies an agent's capabilities as composable, pluggable skills with typed parameters and resource constraints. We formalize the genome space $\genomespace$ and define valid genomes via a constraint satisfaction problem.
    
    \item \textbf{Genome Engine} (\S\ref{sec:genome_engine}): We implement a Genome Engine in Rust that loads, validates, and executes skills dynamically. We present formal algorithms for genome parsing, skill lifecycle management, and hot-reload with consistency guarantees.
    
    \item \textbf{Tiered Memory Architecture} (\S\ref{sec:memory}): We design a three-tier memory system (Hot/Warm/Cold) with tree-structured index achieving O(log n) retrieval, 20-50$\times$ compression through distillation, and bounded context (5KB core) regardless of history size.
    
    \item \textbf{Intelligent Model Router} (\S\ref{sec:router}): We implement a 15-dimension weighted classifier for task complexity, automatic fallback chains, and cost-aware routing that reduces API costs by 10-100$\times$ while maintaining quality.
    
    \item \textbf{Model Health Registry} (\S\ref{sec:health}): We design a circuit-breaker pattern for automatic model failover with persistent state, error classification, and auto-recovery.
    
    \item \textbf{Self-Governance Protocols} (\S\ref{sec:governance}): We implement four protocols (WAL, VBR, ADL, VFM) ensuring autonomous reliability without human supervision.
    
    \item \textbf{Evolutionary Mutation Framework} (\S\ref{sec:evolution}): We define a fitness function $\fitnessfn$ over genome configurations and present mutation operators (skill addition, removal, parameter perturbation, crossover) that enable automated capability search under resource and safety constraints.
    
    \item \textbf{Multi-Tier Deployment Model} (\S\ref{sec:deployment}): We design a deployment architecture spanning bare-metal edge devices (ARMv6, 512\,MB RAM), containerized servers (Podman/Docker), and cloud sandboxes (E2B). Agents synchronize across tiers via MQTT with configurable QoS.
    
    \item \textbf{Production Evaluation} (\S\ref{sec:evaluation}): We evaluate \evoclaw{} on a Raspberry Pi~1, demonstrating 3.0\,MB binary size, 3.2\,MB memory footprint, sub-100\,ms skill load time, and 24+ hour continuous operation with three concurrent skills.
    
    \item \textbf{Open-Source Release}: The complete system---Go orchestrator (5,000 lines), Rust agent (8,000 lines), 383 tests, and documentation---is released under an open-source license for full reproducibility.
\end{enumerate}

\subsection{Paper Organization}

Section~\ref{sec:related} reviews related work. Section~\ref{sec:architecture} presents the system architecture. Section~\ref{sec:genome_engine} details the Genome Engine. Section~\ref{sec:memory} covers tiered memory. Section~\ref{sec:router} describes intelligent routing. Section~\ref{sec:health} covers model health. Section~\ref{sec:governance} presents self-governance protocols. Section~\ref{sec:evolution} describes the evolutionary framework. Section~\ref{sec:deployment} covers multi-tier deployment. Section~\ref{sec:blockchain} discusses blockchain integration. Section~\ref{sec:implementation} discusses implementation details. Section~\ref{sec:evaluation} presents evaluation results. Section~\ref{sec:limitations} discusses limitations. Section~\ref{sec:conclusion} concludes.

%% ========================================================================
%%  SECTION 2: RELATED WORK
%% ========================================================================

\section{Related Work}
\label{sec:related}

\evoclaw{} draws on and extends seven research areas: LLM-based autonomous agents, multi-agent systems, memory architectures, model routing, genetic programming, self-governance, and blockchain for AI.

\subsection{LLM-Based Autonomous Agents}

The emergence of capable LLMs has catalyzed a new wave of autonomous agent research. Wang et al.~\cite{wang2023survey} provide a comprehensive taxonomy of LLM-based agents, identifying four key modules: profiling, memory, planning, and action. Xi et al.~\cite{xi2023rise} survey the landscape from construction to evaluation, emphasizing the gap between current systems and human-level autonomy.

\textbf{Task decomposition agents.} AutoGPT~\cite{auto_gpt} and BabyAGI~\cite{babyagi} pioneered the paradigm of LLM-powered task chaining, where a model decomposes high-level goals into sequences of subtasks. While demonstrating the potential of autonomous planning, these systems cannot modify their own toolset---new capabilities require code changes by human engineers. \evoclaw{} addresses this through genome-driven evolution, enabling autonomous capability acquisition.

\textbf{Lifelong learning agents.} Voyager~\cite{wang2023voyager} represents the closest conceptual precedent to \evoclaw{} in the LLM agent literature. Operating in Minecraft, Voyager maintains an ever-growing skill library of executable JavaScript programs, using GPT-4 to generate new skills via iterative prompting. However, Voyager's skills are confined to a single game environment, are not declaratively specified, and do not support deployment on resource-constrained hardware. \evoclaw{} generalizes this concept to arbitrary domains through declarative skill specifications and multi-tier deployment.

\textbf{Memory-augmented agents.} MemGPT~\cite{packer2023memgpt} draws on operating systems concepts to give LLM agents virtual memory management, enabling extended context through hierarchical memory tiers~\cite{silberschatz2014operating}. While MemGPT uses fixed-size windows with LRU eviction, our tiered architecture uses tree-structured retrieval (inspired by PageIndex~\cite{vectifyai_pageindex}) achieving O(log n) search vs. O(n) vector similarity. Additionally, our distillation engine compresses conversations 20-50$\times$ before storage, addressing MemGPT's context bloat problem.

Reflexion~\cite{shinn2024reflexion} introduces self-reflective agents that critique their own outputs, while Chain-of-Hindsight~\cite{liu2023chain} enables learning from past mistakes. Our WAL (Write-Ahead Log) protocol complements these approaches by ensuring important context survives memory compaction---a reliability pattern not addressed in prior work.

\textbf{Agent frameworks.} LangChain~\cite{langchain} provides composable abstractions for building LLM-powered chains, agents, and tools. CrewAI~\cite{crewai} enables multi-agent collaboration through role assignment. However, both frameworks treat agent architectures as static---changing an agent's capabilities requires code modifications and redeployment. \evoclaw{}'s hot-reloadable genomes eliminate this limitation.

\textbf{Tool use.} ToolFormer~\cite{schick2023toolformer} and Function Calling (OpenAI) enable LLMs to invoke external tools via structured schemas. Our agentic tool loop extends this by implementing multi-turn execution on edge agents, with tool result streaming and automatic error recovery. Unlike centralized tool use~\cite{parisi2023gorilla}, our architecture distributes tool execution across edge, server, and cloud tiers.

\textbf{Software engineering agents.} SWE-agent~\cite{yang2024sweagent} introduces agent-computer interfaces (ACIs) optimized for LLM agents, achieving strong performance on software engineering benchmarks. This work highlights the importance of interface design for agent effectiveness, a principle \evoclaw{} extends to the skill specification layer. Our VBR (Verify Before Reporting) protocol prevents false completion claims---a critical issue for autonomous software agents.

\textbf{Continual learning.} Lamin~\cite{lamar2023lamin} introduces memory-efficient fine-tuning for LLMs. While Lamin focuses on model weights, our evolution operates at the \textit{behavioral} level (skill configurations), enabling adaptation without weight updates. This is more practical for edge deployment where on-device fine-tuning is infeasible.

\subsection{Multi-Agent Systems and Communication}

Multi-agent systems (MAS) have a rich history predating LLMs~\cite{wooldridge2009introduction}. Recent work has focused on LLM-based multi-agent coordination.

\textbf{Role-playing and collaboration.} CAMEL~\cite{li2023camel} proposes communicative agents that collaborate through role-playing, using inception prompting to guide autonomous cooperation. Guo et al.~\cite{guo2024llm_multi_agents} survey LLM-based multi-agent systems across domains, identifying profiling, communication, and capability growth as key challenges. \evoclaw{} addresses capability growth through genome evolution and provides structured communication via MQTT pub/sub.

\textbf{Generative agents.} Park et al.~\cite{park2023generative} create believable simulacra of human behavior using LLM-powered agents with memory, reflection, and planning modules. Their agents exhibit emergent social behaviors in a sandbox environment. Our ADL (Anti-Divergence Limit) protocol prevents the persona drift they observe, ensuring agents maintain consistent behavior over time.

\textbf{Agent evaluation.} AgentBoard~\cite{ma2024agentboard} provides comprehensive benchmarking for multi-turn LLM agents with fine-grained progress metrics. This work highlights the need for rigorous evaluation frameworks, a gap we address with our fitness function formalization and monotonic improvement guarantee.

\textbf{Multi-agent orchestration.} MetaGPT~\cite{ge2023metagpt} assigns different roles (product manager, engineer, architect) to agents for software development. While effective, this requires manual role assignment. Our intelligent router automates task-model assignment via 15-dimension classification, reducing orchestration overhead.

\subsection{Memory Architectures for LLMs}

Memory is critical for LLM agents to maintain context across conversations~\cite{zhong2023token}. Three broad approaches exist:

\textbf{Context window extension.} Methods like Memorizing Contexts~\cite{xiong2023memorizing} train models to remember information within the context window. While effective, this doesn't scale to months-long histories and incurs high inference costs~\cite{tai2022retrieval}.

\textbf{Retrieval-Augmented Generation (RAG).} Dense vector retrieval~\cite{karpukhin2020dense} finds semantically similar documents, then includes them in context~\cite{lewis2020retrieval}. However, similarity $\neq$ relevance: ``Bowen likes coffee'' and ``coffee machine broke'' score high similarity but are unrelated. Our tree-structured retrieval (inspired by PageIndex~\cite{vectifyai_pageindex}) uses LLM reasoning to navigate hierarchies, achieving 98\%+ accuracy vs. 70-80\% for vector RAG.

\textbf{Compressive memory.} Recursively summarizing conversations~\cite{wu2023scaling} reduces context bloat but loses fine-grained details. Our distillation engine achieves 20-50$\times$ compression while preserving critical information through structured fact extraction.

\textbf{Hierarchical memory.} MemGPT~\cite{packer2023memgpt} and HuggingGPT~\cite{shen2023hugginggpt} use tiered memory with different retention policies. While MemGPT uses paging with LRU eviction, our warm tier uses relevance decay scoring (half-life 30 days) and access-count reinforcement---inspired by human memory consolidation~\cite{atkinson1971human}.

\textbf{Neurosymbolic approaches.} KGR-Agents~\cite{wang2024kgr} combine knowledge graphs with vector search for hybrid retrieval. Our tree index is similarly neurosymbolic: LLM reasoning (symbolic) navigates a tree structure (symbolic) to retrieve compressed facts (subsymbolic).

\evoclaw{}'s tiered memory advances the state-of-the-art by: (1) bounding context at 5KB core regardless of history size, (2) achieving O(log n) vs. O(n) retrieval, (3) reasoning-based vs. similarity-based search, and (4) 20-50$\times$ compression through distillation.

\subsection{Model Routing and Orchestration}

With the proliferation of LLMs (hundreds of models across providers), intelligent routing has become critical~\cite{zheng2023survey}.

\textbf{Cost-aware routing.} LLMAdvisor~\cite{liu2023llm} recommends models based on task complexity to reduce costs. FRUGAL~\cite{lyu2023frugal} uses cached small models for routine queries. Our router extends these ideas with 15-dimension weighted classification (vs. keyword heuristics) and automatic fallback chains for reliability.

\textbf{Performance optimization.} DistServe~\cite{zheng2023distserve} and vLLM~\cite{kwon2023efficient} optimize throughput via batching and caching. While orthogonal to our work, our model health registry prevents these systems from routing to degraded models, improving overall cluster reliability.

\textbf{Multi-model ensembles.} Mixture-of-Experts~\cite{shazeer2017outrageously} routes to different model experts based on input. Our router applies this at the \textit{agent} level: simple tasks use budget models, complex tasks use premium models, reasoning tasks use specialized models (e.g., DeepSeek R1).

\textbf{Circuit breaker patterns.} In microservices, circuit breakers~\cite{ford2000circuit} prevent cascading failures by degrading unhealthy services. We adapt this to LLM routing: models with consecutive failures are automatically degraded and requests routed to fallbacks. State persistence across restarts (unlike traditional circuit breakers) improves resilience.

\evoclaw{} advances model routing by: (1) 15-dimension weighted classification (vs. simple heuristics), (2) 5-tier complexity model (Simple/Medium/Complex/Reasoning/Critical), (3) automatic fallback chains (2-3 models per tier), (4) circuit-breaker reliability, and (5) 10-100$\times$ cost reduction while maintaining quality.

\subsection{Genetic Programming and Self-Modifying Systems}

\textbf{Classical GP.} Genetic programming~\cite{koza1992gp} evolves computer programs through evolutionary operators: selection, crossover, and mutation. Traditional GP operates on tree-structured program representations, evolving towards a fitness objective~\cite{angeline1994genetic}. \evoclaw{}'s genome mutation operators are inspired by GP but operate on declarative configuration spaces rather than raw program syntax.

\textbf{LLM-guided evolution.} Eureka~\cite{ma2023eureka} uses LLMs to perform evolutionary optimization over reward functions for reinforcement learning, achieving human-level reward design. This demonstrates the viability of LLM-guided evolutionary search in code spaces. \evoclaw{} applies a similar principle to agent configuration spaces.

\textbf{Self-modifying code.} The concept of programs that modify their own source code dates to early computing~\cite{vonneumann1966self}. Quines~\cite{quine_thompson} demonstrate self-replication, while reflective systems~\cite{smith1984reflection} enable programs to inspect and modify their own execution. \evoclaw{} provides a \textit{safe} form of self-modification by constraining mutations to the declarative genome layer, preventing arbitrary code injection.

\textbf{Continual learning.} Lamin~\cite{lamar2023lamin} introduces memory-efficient fine-tuning. Our evolution is complementary: Lamin adapts model weights, while \evoclaw{} adapts behavioral configurations. Together, they enable dual-level adaptation (neural + symbolic).

\textbf{Deliberate search.} Tree of Thoughts~\cite{yao2023tree} and Language Agent Tree Search (LATS)~\cite{zhou2023lats} apply search algorithms to LLM reasoning, enabling more deliberate problem-solving. \evoclaw{}'s evolutionary operators can be viewed as a form of search over agent configuration space.

\textbf{Multi-objective optimization.} Our fitness function (Eq.~\ref{eq:fitness}) combines multiple objectives via weighted summation. Future work could use Pareto-based approaches like NSGA-II~\cite{deb2002nsga2} for richer trade-off analysis.

\subsection{Self-Governance and Reliability}

Autonomous systems need mechanisms for maintaining reliability without human oversight~\cite{rushby1995critical}. Our four protocols address different failure modes:

\textbf{WAL (Write-Ahead Log).} Write-ahead logging~\cite{hewitt1976implementation} ensures durability in database systems by writing changes before applying them. Our WAL protocol prevents context loss during memory compaction---a pattern not previously applied to LLM agents.

\textbf{VBR (Verify Before Reporting).} Assertions~\cite{meyer1997object} verify program invariants. Our VBR protocol applies this to agent task completion: agents must verify outputs (file exists, tests pass, git push succeeded) before claiming completion.

\textbf{ADL (Anti-Divergence Limit).} Reinforcement learning agents can diverge from intended behavior~\cite{arnold2017constrained}. Our ADL protocol quantifies persona drift via anti-pattern detection (sycophancy~\cite{sharma2023sycophancy}, passivity, hedging) and triggers recalibration when divergence exceeds thresholds.

\textbf{VFM (Value-For-Money).} Cost optimization is critical for production systems~\cite{marr2023cost}. Our VFM protocol tracks task costs vs. outcomes, identifying when cheaper models would suffice (e.g., monitoring tasks using Opus vs. Haiku).

\subsection{Blockchain for AI Agents}

Blockchain technology enables trustless coordination among autonomous agents~\cite{swan2018blockchain}. ClawChain~\cite{clawchain}, our agent-native Layer~1 blockchain, builds on Substrate~\cite{wood2019polkadot} and provides:

\textbf{Decentralized identity.} DIDs (Decentralized Identifiers)~\cite{spencer2015decentralized} give agents cryptographically-verifiable identities without central authorities. ClawChain extends this with reputation tracking: agents accumulate on-chain reputation through verified actions.

\textbf{Agent economies.} Agent markets~\cite{goertzel2017agi} enable agents to buy/sell services. Our genome marketplace allows agents to share evolved configurations: successful genomes can be sold, with ClawChain ensuring payments and reputation-weighted discovery.

\textbf{On-chain governance.} DAOs (Decentralized Autonomous Organizations)~\cite{buterin2014ethereum} enable collective decision-making. Our vision extends this to agents: high-reputation agents propose genome mutations, low-reputation agents adopt them, creating a distributed evolution system.

\textbf{Smart contracts for safety.} Formal verification~\cite{miller2016simpl} ensures smart contracts behave as specified. Our safety invariants (max loss, blocked actions) are enforced via ClawChain smart contracts, preventing malicious genomes from propagating.

\subsection{Edge AI and Lightweight Deployment}

\textbf{TinyML.} The TinyML movement~\cite{tinyml} has driven the development of machine learning models optimized for microcontrollers and resource-constrained devices. Frameworks like TensorFlow Lite Micro enable inference on devices with kilobytes of memory. However, these focus on inference rather than agentic behavior.

\textbf{Edge computing.} The proliferation of IoT devices has spurred research in edge computing~\cite{shi2016edge}, where computation is moved closer to data sources. \evoclaw{}'s multi-tier deployment model is designed for exactly this paradigm, enabling intelligent agents on devices as small as a Raspberry Pi~1.

\subsection{Positioning of \evoclaw{}}

Table~\ref{tab:comparison} positions \evoclaw{} relative to existing systems. To our knowledge, no existing system combines: (1)~declarative agent specification via a genome, (2)~hot-reloadable skills with formal consistency guarantees, (3)~evolutionary mutation operators over agent capabilities, (4)~tiered memory with tree-structured retrieval, (5)~intelligent model routing with circuit-breaker reliability, (6)~self-governance protocols for autonomous operation, (7)~multi-tier deployment from 512\,MB edge devices to cloud, and (8)~production-ready implementation with 90\%+ test coverage.

\begin{longtable}{lcccccccc}
\caption{Comparison of \evoclaw{} with representative agent systems. {\color{coregreen}\checkmark}~=~supported, {\color{corered}$\times$}~=~not supported, {\color{coregray}$\sim$}~=~partial.}
\label{tab:comparison}
\\
\toprule
\textbf{System} & \textbf{Declarative} & \textbf{Hot-} & \textbf{Self-} & \textbf{Tiered} & \textbf{Smart} & \textbf{Self-} & \textbf{Edge} & \textbf{Open} \\
 & \textbf{Spec} & \textbf{Reload} & \textbf{Evolving} & \textbf{Memory} & \textbf{Routing} & \textbf{Gov.} & \textbf{Deploy} & \textbf{Source} \\
\midrule
\endfirsthead

\multicolumn{9}{c}%
{{\tablename\ \thetable{} -- continued from previous page}} \\
\toprule
\textbf{System} & \textbf{Declarative} & \textbf{Hot-} & \textbf{Self-} & \textbf{Tiered} & \textbf{Smart} & \textbf{Self-} & \textbf{Edge} & \textbf{Open} \\
 & \textbf{Spec} & \textbf{Reload} & \textbf{Evolving} & \textbf{Memory} & \textbf{Routing} & \textbf{Gov.} & \textbf